\section{Involution des feuillets dans \(4\) représentations}

L'involution \(C^*\mapsto-C^*\) organise quatre réponses distinctes :
\begin{equation}
\label{rm:eq:modular-sheet-parities}
\begin{aligned}
S_{\rm bi}(A,-C^*)&=S_{\rm bi}(A,C^*),\\
\Gamma_{6\to8}(A,-C^*)&=-\Gamma_{6\to8}(A,C^*),\\
(\widehat\mu,\widehat\varepsilon)&\longleftrightarrow
(\widehat\varepsilon,\widehat\mu),\\
\widehat c&\longmapsto\widehat c,
\qquad D_A\longmapsto D_A.
\end{aligned}
\end{equation}
L'information bicouche est paire ; l'orientation aréale est impaire ; les deux
constitutives forment un doublet, et les déterminants restent fixes. Ce tableau
des parités précède toute interprétation comme symétrie physique.

\section{Transformation rationnelle involutive dans l'espace des paramètres CKM}

Soit
\begin{equation}
\label{rm:eq:modular-Mckm}
M_{\rm CKM}=
\begin{pmatrix}
2&-5&28\\
0&7&-25/2\\
1&-20&-2/3
\end{pmatrix},
\qquad
\det M_{\rm CKM}=-\frac{3857}{6}.
\end{equation}
L'inverse existe sur \(\mathbb Q\) et est
\begin{equation}
\label{rm:eq:modular-Mckm-inverse}
M_{\rm CKM}^{-1}=
\begin{pmatrix}
1528/3857&3380/3857&801/3857\\
75/3857&176/3857&-150/3857\\
6/551&-30/551&-12/551
\end{pmatrix}.
\end{equation}
Avec
\begin{equation}
\label{rm:eq:modular-ckm-coordinates}
h=\frac{C^*}{6},
\qquad
\Delta=\frac{180}{\pi}\Delta_4,
\qquad
\boldsymbol\theta=M_{\rm CKM}(A,h,\Delta)^T,
\end{equation}
la coordonnée impaire entre dans la direction
\begin{equation}
\label{rm:eq:modular-ch}
c_h=\begin{pmatrix}-5\\7\\-20\end{pmatrix}.
\end{equation}

Si \(c_A=(2,0,1)^T\) et
\(c_\Delta=(28,-25/2,-2/3)^T\), les trois colonnes
\((c_A,c_h,c_\Delta)\) sont précisément la base transportée par
\(M_{\rm CKM}\). La deuxième ligne de
\cref{rm:eq:modular-Mckm-inverse} est le covecteur \(\ell_h\) qui extrait la
coordonnée impaire :
\begin{equation}
\label{rm:eq:modular-dual-coordinate-check}
\ell_hc_A=0,
\qquad
\ell_hc_h=1,
\qquad
\ell_hc_\Delta=0.
\end{equation}

\begin{theorem}[Réflexion CKM induite par le feuillet]
\label{rm:thm:modular-ckm-reflection}
Soit
\begin{equation}
\label{rm:eq:modular-Dh-lh}
D_h=\operatorname{diag}(1,-1,1),
\qquad
\ell_h=\frac1{3857}\begin{pmatrix}75&176&-150\end{pmatrix}.
\end{equation}
L'involution induite dans l'espace de \(\boldsymbol\theta\) est
\begin{equation}
\label{rm:eq:modular-Rckm}
\boxed{
\mathcal R_{\rm CKM}
=M_{\rm CKM}D_hM_{\rm CKM}^{-1}
=I-2c_h\ell_h.}
\end{equation}
Satisfait
\begin{equation}
\label{rm:eq:modular-Rckm-properties}
\mathcal R_{\rm CKM}^2=I,
\qquad
\det\mathcal R_{\rm CKM}=-1,
\qquad
\mathcal R_{\rm CKM}M_{\rm CKM}=M_{\rm CKM}D_h.
\end{equation}
\end{theorem}

\begin{proof}
La multiplication de
\cref{rm:eq:modular-Mckm,rm:eq:modular-Mckm-inverse} donne l'identité ; en
particulier, \cref{rm:eq:modular-dual-coordinate-check} prouve que
\(c_h\ell_h\) est le projecteur de rang un sur \(\mathbb Qc_h\) le
long de \(\operatorname{span}\{c_A,c_\Delta\}\). Le produit exact
\(\ell_hc_h=1\) implique
\[
(I-2c_h\ell_h)^2
=I-4c_h\ell_h+4c_h(\ell_hc_h)\ell_h=I.
\]
L'opérateur fixe point par point l'hyperplan
\(\ker\ell_h=\operatorname{span}\{c_A,c_\Delta\}\) et change le signe de
\(c_h\). Ses valeurs propres sont donc \(1,1,-1\) ; d'où
\(\det\mathcal R_{\rm CKM}=-1\) et
\(\Tr\mathcal R_{\rm CKM}=1\). Enfin,
\[
\mathcal R_{\rm CKM}M_{\rm CKM}
=M_{\rm CKM}D_hM_{\rm CKM}^{-1}M_{\rm CKM}
=M_{\rm CKM}D_h,
\]
ce qui prouve l'identité d'entrelacement.
\end{proof}

La matrice explicite est
\begin{equation}
\label{rm:eq:modular-Rckm-matrix}
\mathcal R_{\rm CKM}=
\begin{pmatrix}
4607/3857&1760/3857&-1500/3857\\
-150/551&199/551&300/551\\
3000/3857&7040/3857&-2143/3857
\end{pmatrix}.
\end{equation}
Bien que \(\mathcal R_{\rm CKM}\) ne soit pas une réflexion orthogonale pour le
produit euclidien des trois coordonnées angulaires, elle l'est pour la
métrique rationnelle transportée
\begin{equation}
\label{rm:eq:modular-ckm-metric}
G_{\rm CKM}=(M_{\rm CKM}^{-1})^TM_{\rm CKM}^{-1},
\qquad
\mathcal R_{\rm CKM}^TG_{\rm CKM}\mathcal R_{\rm CKM}=G_{\rm CKM}.
\end{equation}
En effet, dans la base \((c_A,c_h,c_\Delta)\), la métrique est l'identité
et l'involution est \(D_h\). C'est pourquoi le terme ``réflexion'' a un
contenu métrique précis et ne décrit pas seulement le fait que
\(\mathcal R_{\rm CKM}^2=I\).
La phase \(\delta_{\rm CP}=9A\) reste fixe sous cette réflexion. C'est
pourquoi \(\mathcal R_{\rm CKM}\) n'est pas la conjugaison CP physique, mais le
transport de l'inversion de feuillet HMT à l'espace de mélange.

L'inverse rationnel de \(M_{\rm CKM}\) retrouve
\begin{align}
A&=\frac{1528\theta_{12}+3380\theta_{23}+801\theta_{13}}{3857},
\label{rm:eq:modular-A-from-ckm}\\
C^*&=\frac{6(75\theta_{12}+176\theta_{23}-150\theta_{13})}{3857}.
\label{rm:eq:modular-C-from-ckm}
\end{align}
La troisième coordonnée se retrouve également sans perte :
\begin{equation}
\label{rm:eq:modular-Delta-from-ckm}
\Delta
=\frac{6\theta_{12}-30\theta_{23}-12\theta_{13}}{551}.
\end{equation}
Ces équations sont des changements de coordonnées. Elles ne font pas de CKM la cause
rétrospective de \(A,C^*\) ou de \(\gammaH\).

Comme reconnaissance terminale, le constructeur préalable détermine
\begin{equation}
\label{rm:eq:modular-ckm-values}
(\theta_{12},\theta_{23},\theta_{13},\delta_{\rm CP})
=(13.003313589509^\circ,
2.398056157332^\circ,
0.213977382244^\circ,
65.676173123554^\circ).
\end{equation}
La réflexion rationnelle est démontrée avant l'utilisation de ces décimales.

\begin{keyresult}
La transition \(6\to8\) est d'abord formalisée par l'inclusion
\(\iota_{6,8}:\mathcal F_6\hookrightarrow\mathcal F_8\), dont on
obtient
\(90,120\) et le défaut \(-1/12\). Le mot de six trits fournit
ensuite la bijection \(729\leftrightarrow729\), et
\(\mathcal J_{6\to8}\) transporte les deux modes collectifs au bord. Sur
cet antécédent, l'aire et le hamiltonien modulaire reconstruisent exactement
les deux familles ; la purification thermochamp double détermine l'entropie
d'intrication, et la première loi
finie conserve le défaut d'entropie relative. Dans CKM, l'inversion de
feuillet devient une réflexion rationnelle de rang impair un.
\end{keyresult}

\begin{falsifier}
La généalogie initiale échoue si \(\iota_{6,8}\) n'est pas injective, si les
décomptes ne sont pas \(90/120\), si \(\delta_{6,8}^{\rm inc}\neq-1/12\), si
\(\beta_{729}\) n'est pas bijective ou si elle est présentée à tort comme un
isomorphisme additif. Le transport collectif échoue si
\(\mathcal J_{6\to8}D_{\rm inc}\neq
D_{\rm area}\mathcal J_{6\to8}\). La reconstruction échoue si le
déterminant des observables de bord n'est pas trente ou si
\cref{rm:eq:modular-inverse-channels} ne retrouve pas les occupations.
La TFD échoue si sa trace partielle n'est pas \(\rho_A\). La première loi ne
s'étend pas aux changements finis sans l'entropie relative de
\cref{rm:eq:modular-finite-first-law}. Le type
\(III_{\lambda_\partial}\) est invalidé si un canal ne persiste pas, s'il n'y a pas de
produit compatible ou si le groupe des rapports change. La réflexion CKM
est réfutée par n'importe laquelle des conditions
\(\mathcal R^2\neq I\), \(\det\mathcal R\neq-1\) ou
\(\mathcal RM\neq MD_h\). Une formule
\(\boldsymbol\theta=f(\gammaH)\) sans entrelaceur viole l'indépendance
typée des deux réalisations fonctionnelles de l'antécédent commun.
\end{falsifier}

\begin{statusbox}
Inclusion \(H\subset O\), dénombrements \(90/120\), deux normalisations du
défaut \(-1/12\), bijection de six trits \(729\leftrightarrow729\),
entrelaceur collectif, échelle \(a_0\), opérateurs finis, spectre d'aire,
TFD, défaut orienté, première loi locale et finie et réflexion CKM :
\status{Théorème interne exact avec ses structures initiales explicites}.
Classification
\(III_{\lambda_\partial}\) : \status{Théorème avec hypothèses explicites de
persistance}. Interprétation de la copie thermochamp double comme extérieur physique :
\status{Réalisation physique conditionnelle}. La grandeur \(\gammaH\) est reprise
du chapitre consacré à Barbero--Immirzi et à l'aire minimale, où elle a été dérivée ; elle intervient ici
comme donnée structurelle déjà établie.
\end{statusbox}
